% Supported v0.5 research-statement example.
%
% Compile with LuaLaTeX:
%   latexmk -lualatex -interaction=nonstopmode -halt-on-error \
%     examples/statements/research-statement.tex
\documentclass[type=research]{careerdossier-statement}
% examples/profiles/profile-academic.tex
%
% Shared profile metadata used by the academic-CV example and the corresponding
% industry-resume smoke fixture. Keeping it independent of document layout
% proves that both document families consume the same \CDossierSetup interface.

\CDossierSetup{
  name     = {Ada Lovelace},
  headline = {Researcher in Analytical Computing},
  email    = {ada.lovelace@example.com},
  location = {Toronto, Ontario, Canada},
  website  = {example.com/ada-lovelace},
  %% The bare Google Scholar identifier, expanded to
  %% scholar.google.com/citations?user=ada-example for both the displayed text
  %% and the link. Pasting the whole address works too and is left as written.
  %% See docs/API.md, "Accepted forms for the web-profile keys".
  scholar  = {ada-example},
  orcid    = {0000-0002-1825-0097},
  affiliation = {Example University},
}

\CDossierStatementSetup{
  subtitle={Transparent methods for trustworthy scientific computing},
  application-context={Application for Assistant Professor of Computational Science},
  application-id={APP-2026-0042}
}
\begin{document}
\MakeCDossierStatementHeader

\section*{Research Vision and Significance}

My research asks how computational results can remain understandable,
reproducible, and useful after the immediate project that produced them has
ended. I develop numerical methods and research-software practices that expose
their assumptions, preserve the path from evidence to conclusion, and allow
scientists from neighbouring fields to test and extend the work. The central
premise is that reliability is not a final verification step: it is an
organizing principle for model design, implementation, documentation, and
communication.

This programme connects mathematical analysis with the practical conditions in
which scientific software is built. A method is most valuable when its
guarantees, software realization, and limits can be inspected together. That
connection matters in collaborative research, where a result may pass through
several disciplines and computing environments before it becomes evidence for a
decision.

\section*{Prior and Current Contributions}

My first line of work develops adaptive algorithms whose error estimates remain
meaningful when data quality and computational resources vary. I combine
analysis with executable experiments, using small reference implementations to
make each theoretical claim testable. This approach helps collaborators
distinguish uncertainty arising from a physical model from error introduced by
discretization, data preparation, or implementation choices.

A second line of work concerns durable computational records. I design workflows
that connect source data, environment descriptions, tests, and figures without
requiring readers to reconstruct an undocumented sequence of manual steps. In
collaborative projects, these records support method comparison, onboarding, and
independent replication. They also give students a concrete model of research
accountability: a result is incomplete until another person can understand how
it was produced.

These projects share a commitment to accessibility across expertise. I write
software interfaces and technical explanations for readers who know the
scientific question but may not know the implementation language or numerical
method. This requires choosing meaningful defaults, naming uncertainty clearly,
and separating essential decisions from incidental technical detail.

\section*{Future Research Programme}

The next phase has three connected aims. First, I will develop verification
methods for computational pipelines that combine learned and mechanistic
models. The goal is to identify which guarantees can be preserved, which
assumptions must be evaluated empirically, and how uncertainty should be
reported when components change over time.

Second, I will study reproducibility as a property of research infrastructure
rather than an after-the-fact archive. This work will evaluate lightweight ways
to record environments, data transformations, and validation decisions during
ordinary scientific practice. Expected outputs include open reference
implementations, reusable evaluation datasets, and guidance that small research
groups can adopt without a dedicated engineering team.

Third, I will build teaching and mentorship into the programme. Students will
participate in paired analysis and implementation reviews, learn to design tests
from mathematical claims, and publish documentation for readers outside the
immediate project. This structure gives students ownership of substantive
questions while making standards for reliable work explicit.

\section*{Institutional Fit and Trajectory}

Example University's strengths in computational science, applied mathematics,
and open research create a productive setting for this agenda. I would seek
collaborations with researchers whose domain questions stress current methods,
and I would contribute shared software and training that make those
collaborations easier to sustain. Over the next five years, I aim to establish a
research group known for connecting rigorous numerical reasoning with tools and
practices that other scientists can genuinely reuse.

\end{document}
